\section{NLP experiments details}
\label{sec:apndx-expt-nlp}

\subsection{MLM Pretraining}
\label{sec:app-expt-nlp:mlm}

We use four publicly available datasets Books \citep{zhu2015aligning}, CC-News \citep{guu2020realm}, Stories \citep{trinh2018simple} and Wikipedia to pretrain \bigb. 
We borrow the sentencepiece vocabulary as RoBERTa (which is in turn borrowed from GPT2).
We split any document longer than $4096$ into multiple documents and we join documents that were much smaller than $4096$.
Following the original BERT training, we mask $15\%$ of tokens in these four datasets, and train to predict the mask. We warm start from RoBERTa's checkpoint. 
We train two different models: \bigb-\textsc{itc}-base and \bigb-\textsc{etc}-base. The hyper-parameters for these two models are given in~\Cref{tab:app_mlm_param}. In all experiments we use a learning
rate warmup over the first 10,000 steps, and linear
decay of the learning rate. 

Similar to the norm, we trained a large version of model as well, which has 24 layers with 16 heads and hidden dimension of 1024.
Following the observation from RoBERTa, we pretrain on a larger batch size of 2048 for this size.
For \bigb-\textsc{itc} the block length was kept same as base size, but for \bigb-\textsc{etc} the block length was almost doubled to 169. All the remaining parameters were the same.


\begin{table}[bht]
\small
\centering
\begin{tabular}{@{}l c r c r @{}}
\toprule
Parameter & & \bigb-\textsc{itc} & & \bigb-\textsc{etc} \\
\midrule
 Block length, $b$ & & $64$ & & 84\\
 $\#$ of global token, $g$ & & $2\times b$ & & $256$ \\
 Window length, $w$  & & $3\times b$ & & $3\times b$ \\
 $\#$ of random token, $r$ & & $3\times b$ & & $0$ \\
 Max. sequence length & & $4096$ & & $4096$\\
 $\#$ of heads & & $12$ & & $12$\\
 $\#$ of hidden layers & & $12$ & & $12$ \\
 Hidden layer size & & $768$ & & $768$ \\
 Batch size & & $256$ & & $256$ \\
 Loss & & MLM & & MLM \\
 Activation layer & & gelu & & gelu \\ 
 Dropout prob & & $0.1$ & & $0.1$ \\
 Attention dropout prob & & $0.1$ & & $0.1$ \\
 Optimizer & & Adam & & Adam\\
 Learning rate & & $10^{-4}$  & & $10^{-4}$\\
 Compute resources & & $8 \times 8$ TPUv3 & & $8 \times 8$ TPUv3\\
\bottomrule
\end{tabular}
\vspace{2mm}
\caption{Hyperparameters for the two \bigb base models for MLM.}
\label{tab:app_mlm_param}
\end{table}

\begin{table}[b]
    \centering
    \small
    \parbox{.45\linewidth}{
    \centering
    \begin{tabular}{@{}lrr@{}}
    \toprule
    Dataset & $\#$ tokens & Avg. doc len. \\
    \midrule
        Books \citep{zhu2015aligning} & $1.0$B & $37$K \\
        CC-News \citep{guu2020realm} & $7.4$B & $561$ \\
        Stories \citep{trinh2018simple} & $7.7$B & $8.2$K \\
        Wikipedia &  $3.1$B & $592$ \\
    \bottomrule
    \end{tabular}
    \vspace{2mm}
    \caption{Dataset used for pre training.}
    \label{tab:mlm_data}
    }
    \quad \quad
    \parbox{.45\linewidth}{
    \centering
    \begin{tabular}{@{}lcc@{}}
    \toprule
    Model &  Base & Large \\
    \midrule
    RoBERTa (sqln: 512)  & 1.846 & 1.496 \\
    Longformer (sqln: 4096)  & 1.705 & 1.358 \\
    \bigb-\textsc{itc} (sqln: 4096)   & 1.678 & 1.456 \\
    \bigb-\textsc{etc} (sqln: 4096)   & \textbf{1.611} & \textbf{1.274} \\
     \bottomrule
    \end{tabular}
    \vspace{2mm}
    \caption{MLM performance on held-out set.}
    \label{tab:mlm_bpc}
    }
\end{table}

\begin{table}
    \small
    \centering
    \begin{tabular}{@{}lrrrrrr@{}}
    \toprule
    & & \multicolumn{2}{c}{Instances} & & \multicolumn{2}{c}{Instance Length} \\
    \cmidrule{3-4} \cmidrule{6-7}
    Dataset & & Training & Dev & & Median & Max \\
    \midrule
    HotpotQA-distractor \citep{yang2018hotpotqa} & & $90447$ & $7405$ & & $1227$ & $3560$ \\
        Natural Questions \citep{kwiatkowski2019natural} & & $307373$ & $7830$ & & $3258$ & $77962$ \\
        TriviaQA \citep{JoshiTriviaQA2017} & & $61888$ & $7993$ & & 4900 & 32755 \\
        WikiHop \cite{welbl2018constructing} & & $43738$ & $5129$ & & $1541$ & $20337$ \\
    \bottomrule
    \end{tabular}
    \vspace{2mm}
        \caption{Question Answering Datasets}
    \label{tab:qa_data}
\end{table}

\begin{table}[t]
\small
\centering
\begin{tabular}{@{}l c rr c rr c rr c rr@{}}
\toprule
Parameter & & \multicolumn{2}{c}{HotpotQA} & & \multicolumn{2}{c}{NaturalQ} & & \multicolumn{2}{@{}c@{}}{TriviaQA} & & \multicolumn{2}{@{}c@{}}{WikiHop}\\
\cmidrule{3-4} \cmidrule{6-7} \cmidrule{9-10} \cmidrule{12-13}
Global token location & & \textsc{itc} &\textsc{etc} & & \textsc{itc} & \textsc{etc} & &\textsc{itc} & \textsc{etc}& &\textsc{itc} & \textsc{etc} \\
\midrule
$\#$ of global token, $g$ & & $128$ & $256$ & & $128$  & $230$ & & $128$  & $320$ & & $128$ & $430$   \\
 Window length, $w$  & & $192$ & $252$ & & $192$ & $252$ & &  $192$ & $252$ & & $192$ & $252$  \\
 $\#$ of random token, $r$ & & $192$ &  $0$ & & $192$ &  $0$ & & $192$ &  $0$ & & $192$ &  $0$ \\
 Max. sequence length & & $4096$ & $4096$ & & $4096$ & $4096$ & & $4096$ &  $4096$ & & $4096$ &  $4096$\\
 $\#$ of heads & & $12$ & $12$ & & $12$ & $12$ & & $12$ & $12$ & & $12$ & $12$\\
%   Base / Large & & Large & & Large & & Large \\
 $\#$ of hidden layers & & $12$ & $12$ & & $12$ & $12$ & & $12$ & $12$ & & $12$ & $12$ \\
 Hidden layer size & & $768$ & $768$ & & $768$ & $768$ & & $768$ & $768$ & & $768$ &   $768$ \\
 Batch size & & $32$ & $32$ & & $128$ & $128$ & & $32$ &  $32$ & & $64$ &  $64$ \\
 \multirow{2}{*}{Loss} & &   \multicolumn{2}{c}{cross-entropy} & & \multicolumn{2}{c}{cross-entropy} & & \multicolumn{2}{@{}c@{}}{cross-entropy} & & \multicolumn{2}{@{}c@{}}{cross-entropy} \\
 & &   \multicolumn{2}{c}{golden spans} & &   \multicolumn{2}{c}{golden spans} & & \multicolumn{2}{@{}c@{}}{noisy spans~\citep{clark2017simple}} & & \multicolumn{2}{@{}c@{}}{ans choices} \\
 Compute resources & & \multicolumn{2}{c}{$4 \times 2$ TPUv3} & & \multicolumn{2}{c}{$4 \times 8$ TPUv3} & &  \multicolumn{2}{@{}c@{}}{$4 \times 2$ TPUv3} & &  \multicolumn{2}{@{}c@{}}{$4 \times 4$ TPUv3}\\
\bottomrule
\end{tabular}
\vspace{2mm}
\caption{Hyperparameters of base \bigb model used for Question Answering i.e.~the numbers reported in \Cref{tab:QADev}}
\label{tab:app_qa_dev}
\end{table}


\begin{table}[t]
\small
\centering
\begin{tabular}{@{}l c r c r c r c r@{}}
\toprule
Parameter & &  HotpotQA & & NaturalQ & & TriviaQA & & WikiHop \\
\midrule
Global token location & & \textsc{etc} & & \textsc{etc} & & \textsc{etc} & &\textsc{etc} \\
 $\#$ of global token, $g$ & & $256$ & & $230$ & & $320$ & & $430$  \\
 Window length, $w$  & & $507$ & & $507$ & & $507$ & & $507 $ \\
 $\#$ of random token, $r$ & & $0$ & & $0$ & & $0$ & & $0$ \\
 Max. sequence length & & $4096$ & & $4096$ & & $4096$ & & $4096$\\
 $\#$ of heads & & $16$ & & $16$ & & $16$ & & $16$\\
%   Base / Large & & Large & & Large & & Large \\
 $\#$ of hidden layers & & $24$ & & $24$ & & $24$ & & $24$ \\
 Hidden layer size & & $1024$ & & $1024$ & & $1024$ & & $1024$ \\
 Batch size & & $32$ & & $64$ & & $32$ & & $64$ \\
 Loss & & cross-entropy & & cross-entropy & & cross-entropy & & cross-entropy \\
 Num epochs & & $\{5, 9\}$ & & $\{3,5\}$ & & $\{3,5\}$ & & $\{5, 10\}$ \\
 Optimizer  & & Adam  & & Adam & & Adam & & LAMB\\
 Learning rate  & & $3\times 10^{-5}$ & & $\{5, 10\}\times 10^{-5}$ & & $\{3, 5\}\times 10^{-5}$ & & $\{2,5 \}\times 10^{-5}$ \\
 Compute resources & & $4 \times 4$ TPUv3 & & $4 \times 8$ TPUv3 & & $4 \times 4$ TPUv3 & & $4 \times 8$ TPUv3\\
\bottomrule
\end{tabular}
\vspace{2mm}
\caption{Hyperparameters of large \bigb model for Question Answering submitted for test i.e.~the numbers reported in~\Cref{tab:QATest}}
\label{tab:app_qa}
\end{table}

\subsection{Question Answering}
\label{sec:app-expt-nlp:qa}

The detailed statistics of the four datasets used are given in~\Cref{tab:qa_data}. 
All the hyperparameters for \bigb, used for creating~\Cref{tab:QADev} are shown in~\Cref{tab:app_qa_dev} and those submitted to get~\Cref{tab:QATest} 
are shown in \Cref{tab:app_qa}. 
We use two types of regularization in training:

\begin{itemize}[leftmargin=12mm, itemsep=0mm, partopsep=0pt,parsep=0pt]
    \item We used a variant of contrastive predictive coding~\citep{oord2018representation} as a dual encoder model.
    \item We use position embedding for \textsc{itc} and relative position encoding~\citep{shaw2018self} for \textsc{etc}.
\end{itemize}


Next, we will mention the dataset/task specific part of the model.

\paragraph{HotpotQA}
The data consists of each question with multiple evidence paragraphs. 
We filtered 16 QA where the answer was not in the given  evidences. 
For \bigb-\textsc{itc}, we use first $128$ global tokens. 
For \bigb-\textsc{etc}, we have one global token for each question token, one for each evidence paragraph, and one for each sentence within the paragraph, for a maximum of $256$ global token.
We use a dense layer on the output corresponding to global token of the evidence paragraph to predict whether its a supporting fact with a threshold over the output logits.
The answer type (yes/no/span) is predicted with a single dense layer from the global CLS token.
For span based answers, the spans are predicted with dense layers on the sequence with the distance between start and end positions to be no more than 30 words. 
The spans are ranked by sum of start and end logits.


\paragraph{Natural Questions}
Here also the data consists of question with supporting evidence, but in form of a single, potentially long, document and not multiple paragraphs. 
We largely follow the setup of \citep{alberti2019bert}.
For documents, that are longer than
4096, a sliding window approach is used
with stride of 2048.
We use CLS token at the beginning, followed by the question followed by a separator token followed by the document as input.
For \bigb-\textsc{itc}, we make the first $128$ tokens as global. For \bigb-\textsc{etc}, we make a global token for CLS, question, and one token for each of the paragraphs.
We train four predictors at the final layer to predict long answer start, long answer end, short answer start and short answer end respectively.
Instead of independently predicting the start and end of answers we first predict the start and then predict the best end location beyond the start.
For short answer, we limit the distance between start and end positions to be no more than 38 words.
The answer type (null, yes, no, short, long) is predicted from CLS token output embedding.
When the logit for a yes/no answer is higher than the logits for short, long or null answer, we replace the short answer with a corresponding yes/no text.

\paragraph{TriviaQA}
The data consists of question-answer pairs with Wikipedia articles as the ``noisy'' supporting evidence.
We call them noisy because the given Wikipedia articles may or may not contain the answer.
Moreover, the answer entities is not annotated to appropriate span in the article, rather all occurrences found using fuzzy string matching are listed.
We use CLS token at the beginning, followed by the question followed by a separator token followed by the document as input.
For \bigb-\textsc{itc}, we make the first $128$ tokens as global. 
For \bigb-\textsc{etc}, we make a global token for CLS, question, and one token for each sentence up to a maximum of 320 global tokens.
Given the noisy nature of answer span, we follow~\citet{clark2017simple} for training.
We use a dense layer on the sequence to predict the answer span for each article independently, with the distance between start and end positions to be no more than 16 words.
For each article the span with maximum start logit + end logit is chosen.
Then we normalize over all the documents associated with that question.


\paragraph{WikiHop}
For each question in WikiHop, we are given upto $79$ candidates, and $63$ supporting paragraphs.
In our \bigb-\textsc{itc} model, following~\citet{beltagy2020longformer}, we concatenate the answer and the question with special tokens, 
\texttt{[q] Question [/q] [ans] Ans1 [/ans] $\ldots$ [ans] AnsN [/ans]} along with the context. 
As the start of the text, always contains questions followed by answers, we make the first $128$ token attend globally. 
In \bigb-\textsc{etc} model, we do not need to insert special \texttt{[ans]}, \texttt{[/ans]} etc.~as we design global tokens appropriately.
Along with global tokens for question, we have one per candidate answer up to a maximum of 430.
Further, we linked answer tokens to their mentions using relative position label. 
Lastly, we use a dense layer that takes in the output vector corresponding to a candidate answer, and predicts a score for the current candidate to be the correct answer.
We apply this dense layer to each candidate independently and 
the candidate with the best score is picked as our final answer.


It is worthwhile to note that explicitly designed attention connection in \textsc{etc} works slightly better, the random connection based \textsc{itc} is pretty competative.

\subsection{Relationship to Contemporary Work} \label{sec:app-related-work}

\paragraph{Longformer}
\citet{child2019generating} introduced localized sliding window to reduce computation.
A more recent version, which includes localized sliding windows and global tokens was introduced independently by
Longofrmer\citep{beltagy2020longformer}. Although \bigb contains additional random tokens, there are also differences in the way global and local tokens are realized. In particular even when there is no random token, as used to get SoTA in question answering, there are two key differences between Longformer and \bigb-etc (see \citep{ainslie2020etc}):
\begin{enumerate}
    \item We use global-local attention with relative
position encodings enables it to better handle structured inputs
    \item Unlike Longformer, we train the global tokens using CPC loss and learn their use during finetuning.
\end{enumerate}

\subsection{Classification}
\label{sec:app-expt-nlp:cls}
We try two types of classification task.

\begin{table}[bht]
\small
\centering
\begin{tabular}{@{}l c r r r r r @{}}
\toprule
Parameter & & IMDb & Arxiv & Patents & Hyperpartisan & Yelp-5  \\
\midrule
 Batch size & & 64 & 64 & 64 & 32 & 32 \\
 Learning rate & & $1\times 10^{-5}$ & $3\times 10^{-5}$ & $5\times 10^{-5}$ & $5\times 10^{-6}$ & $2\times 10^{-5}$\\
 Num epochs & & 40 & 10 & 3 & 15 & 2 \\
 TPUv3 slice  & & $4 \times 4$ & $4 \times 4$ & $4 \times 4$ & $4 \times 2$ & $4 \times 8$ \\
 $\#$ of heads & & \multicolumn{4}{c}{12} & 16\\
 $\#$ of hidden layers & & \multicolumn{4}{c}{12} & 24 \\
 Hidden layer size & & \multicolumn{4}{c}{768} & $1024$ \\
 Block length, $b$ & & \multicolumn{5}{c}{64} \\
 Global token location & & \multicolumn{5}{c}{\textsc{itc}} \\
 $\#$ of global token, $g$ & & \multicolumn{5}{c}{$2\times b$} \\
 Window length, $w$  & & \multicolumn{5}{c}{$3\times b$ } \\
 $\#$ of random token, $r$ & & \multicolumn{5}{c}{$3\times b$} \\
 Max. sequence length & & \multicolumn{5}{c}{4096}\\
 Vocab size & & \multicolumn{5}{c}{$50358$}\\
 Activation layer & & \multicolumn{5}{c}{gelu} \\ 
 Dropout prob & & \multicolumn{5}{c}{0.1} \\
 Attention dropout prob & & \multicolumn{5}{c}{0.1} \\
 Loss & & \multicolumn{5}{c}{cross-entropy} \\
 Optimizer & & \multicolumn{5}{c}{Adam}\\
\bottomrule
\end{tabular}
\vspace{2mm}
\caption{Hyperparameters for document classification.}
\label{tab:app_dc}
\end{table}

\begin{table}
    \centering
    \small
    \begin{tabular}{@{}lrrrrr@{}}
    \toprule
    Model & IMDb~\citep{maas2011learning} & Yelp-5~\citep{zhang2015character} & Arxiv~\citep{he2019long}  & Patents~\citep{lee2020patent} & Hyperpartisan~\citep{kiesel2019semeval}\\
    \midrule
    \# Examples & 25000 & 650000 & 30043 & 1890093 & 645 \\
    \# Classes & 2 & 5 & 11 & 663 & 2 \\
    Excess fraction & 0.14 & 0.04 & 1.00 & 0.90 & 0.53 \\
    \midrule
    SoTA & \citep{thongtan2019sentiment} 97.4 & \citep{abreu2019hierarchical} 73.28 & \citep{olson2019adapting} 87.96 &  \citep{olson2019adapting} 69.01 &
    \citep{jiang2019team} 90.6 \\
    RoBERTa & $95.0 \pm 0.2$ & 71.75 & 87.42 &  67.07 & $87.8 \pm 0.8$ \\
    \bigb  & $95.2\pm0.2$ & 72.16 & \textbf{92.31} & 69.30 & $\mathbf{92.2 \pm 1.7}$ \\
     \bottomrule
    \end{tabular}
    \vspace{2mm}
        \caption{Classification results. We report the F1 micro-averaged score for all datasets. Experiments on smaller IMDb and Hyperpartisan datasets are repeated 5 times and the average performance is presented along with standard deviation.}
    \label{tab:cls}
\end{table}

\begin{table}[bht]
\small
\centering
 \begin{tabular}{@{}lcccccccc@{}}
    \toprule
System             &  MNLI-(m/mm) & QQP  & QNLI & SST-2 & CoLA & STS-B & MRPC & RTE  \\
                   & 392k         & 363k & 108k & 67k   & 8.5k & 5.7k  & 3.5k & 2.5k   \\ 
\midrule
BERT               & 84.6/83.4    & 71.2 & 90.5 & 93.5  & 52.1 & 85.8  & 88.9 & 66.4 \\
XLNet              & 86.8/-       & 91.4 & 91.7 & 94.7  & 60.2 & 89.5  & 88.2 & 74.0 \\
RoBERTa            & 87.6/-       & 91.9 & 92.8 & 94.8  & 63.6 & 91.2  & 90.2 & 78.7 \\
\bigb     & 87.5/87.3    & 88.6 & 92.2 & 94.6  & 58.5 & 87.8  & 91.5 & 75.0 \\
    \bottomrule
   \end{tabular}
   \vspace{2mm}
      \caption{GLUE Dev results on base sized models. 
   Number of training examples is reported below each task.
   MCC score is reported for CoLA, F1 score is reported for MRPC, Spearman correlation is reported for STS-B, and accuracy scores are reported for the other tasks.}
   \label{tab:glue_dev}
\end{table}

\paragraph{Document classification}
We experiment on datasets of different lengths and contents,
as listed in ~\Cref{tab:cls}.
In particular, we look at sentiment analysis (IMDb~\citep{maas2011learning} and Yelp-5~\citep{zhang2015character}) task and topic assignment (Arxiv~\citep{he2019long}, Patents~\citep{lee2020patent}, and Hyperpartisan~\citep{kiesel2019semeval}) task.
Following BERT, we used one layer with cross entropy loss on top of the first [CLS] token from the \bigb encoder consuming 4096 tokens.
We report the results of document classification experiments in~\Cref{tab:cls}.
We compare against state-of-the-art (SoTA) methods for each dataset and plain RoBERTa model with 512 tokens truncation.
In all experiments we use a learning rate warmup over the first 10\% steps, and linear decay of the learning rate and detail list of remaining hyperparameters are provided in
\Cref{tab:app_dc}.
For better quantitative evaluation, we compute the fraction of the dataset that exceeds 512 tokens, i.e. the length at which the document are often truncated.
We see that gains of using \bigb are more significant when we have longer documents and fewer training examples.
For instance, using base sized model,
\bigb improves state-of-the-art for Arxiv dataset by about $\bm{5\%}$ \textbf{points}.
On Patents dataset, there is improvement over using simple BERT/RoBERTa, but given the large size of training data the improvement over SoTA (which is not BERT based) is not significant.
Note that this performance gain is not seen for much
smaller IMDb dataset.
Along with experimental setup detail, we present detailed results in~\Cref{sec:app-expt-nlp:cls} which show competitive performance.



\paragraph{GLUE}
The General Language Understanding Evaluation (GLUE) benchmark \citep{wang2018glue}, test language models on 8 different natural language understanding tasks.
We used the same training parameters as mentioned in \url{https://github.com/pytorch/fairseq/blob/master/examples/roberta/README.glue.md}. Our model parameters are $b=64, g=2\times b, w = 3 \times b, r = 3 \times b$ ( we used the \bigb-\textsc{itc} base model pretrained on MLM task). 
We compare the performance of \bigb to BERT, XLNet \citep{yang2019xlnet} and RoBERTa in \Cref{tab:glue_dev}. We find that even on task that have a much smaller context, our performance is competitive to full attention models.


\subsection{Summarization\label{sec:appn_summarization}}
As discussed in~\Cref{sec:seq2seq},
given the small length of output sequence, we used sparse \bigb attention only for encoder, while keeping the full attention for decoder.
The number of hidden layers, number of heads, and hidden dimension is same for encoder and decoder.
The hyperparameters are detailed in \Cref{tab:app_sum}.
We summarize our result in~\Cref{tab:app_sum_num}. In all experiments, we use a learning
rate warmup over the first 10,000 steps, and square root
decay of the learning rate.

\begin{table}[bht]
\small
\centering

\begin{tabular}{@{}l r r c r @{}}

\toprule
Parameter & \multicolumn{2}{r}{Base: \bigb-RoBERTa} & & Large: \bigb-Pegasus \\
\midrule

 Block length, $b$ & & $64$ & & $64$\\
 Global token location & & \textsc{itc} & & \textsc{itc} \\
 $\#$ of global token, $g$ & & $2\times b$ & & $2\times b$ \\
 Window length, $w$  & & $3\times b$ & & $3\times b$ \\
 $\#$ of random token, $r$ & & $3\times b$ & & $3\times b$ \\
 \multirow{3}{*}{Max. encoder sequence length} & & BBC-XSUM: \qquad \phantom{.}1024 & & 1024 \\
 & & CNN/DM: \qquad \phantom{.}2048  & & 2048 \\
 & & Others: \qquad \phantom{.}3072 & & 3072 \\
 \multirow{3}{*}{Max. decoder sequence length} & & BBC-XSUM: \qquad \phantom{10.}64 & & 64 \\
 & & CNN/DM: \qquad \phantom{1.}128  & & 128 \\
 & & Others: \qquad \phantom{1.}256  & & 256 \\
 Beam size & & 5 & & 5 \\
 \multirow{2}{*}{Length penalty} & & BBC-XSUM: \qquad \phantom{10}0.7 & & 0.7 \\
 & & Others: \qquad \phantom{10}0.8  & & 0.8 \\
 $\#$ of heads & & $12$ & & $16$\\
 $\#$ of hidden layers & & $12$ & & $16$ \\
 Hidden layer size & & $768$ & & $1024$ \\
 Batch size & & $128$ & & $128$ \\
 \multirow{2}{*}{Loss} & & teacher forced & & teacher forced \\
 & & cross-entropy & & cross-entropy \\
 Activation layer & & gelu & & gelu \\ 
 Dropout prob & & $0.1$ & & $0.1$ \\
 Attention dropout prob & & $0.1$ & & $0.1$ \\
 Optimizer & & Adam & & Adafactor\\
 Learning rate & & $1\times10^{-5}$  & & $1\times10^{-4}$\\
  Compute resources & & $4 \times 4$ TPUv3 & & $4 \times 8$ TPUv3\\
\bottomrule
\end{tabular}
\vspace{2mm}
\caption{Encoder hyperparameters for Summarization. We use full attention in decoder}
\label{tab:app_sum}
\end{table}

\begin{table}[bht]
    \centering
    \small
    \begin{tabular}{@{}lrrrrrccrcc@{}}
    \toprule
    & & \multicolumn{3}{c}{Instances} & & \multicolumn{2}{c}{Input Length}& & \multicolumn{2}{c}{Output Length} \\
    \cmidrule{3-5} \cmidrule{7-8} \cmidrule{10-11}
    Dataset & & Training & Dev & Test & & Median & 90\%-ile & &  Median & 90\%-ile \\
    \midrule
       Arxiv \citep{cohan2018discourse} & & 203037 & 6436 & 6440 & & 6151 & 14405 & & 171 & 352 \\
       PubMed \citep{cohan2018discourse} & & 119924 & 6633 & 6658 & & 2715 & 6101 & & 212 & 318 \\
       BigPatent \citep{sharma2019bigpatent} & &  1207222 & 67068 & 67072 & & 3082 & 7693  & & 123 & 197  \\
    \bottomrule
    \end{tabular}
    \vspace{2mm}
     \caption{Statistics of datasets used for summarization.}
    \label{tab:long_sum_data}
\end{table}

Following success of several recent works~\citep{rothe2019leveraging,liu2019roberta}, we warm start our encoder-decoder \bigb transformer model with pretrained weights and the weights between encoder and decoder are shared.
In particular, the query/key/value matrix of self-attention and all the feedforward layers are shared between encoder and decoder.
The only variable that is initialized randomly is the encoder-decoder attention.
For base sized model, we utilize our MLM pretrained model on 4096 sequence length from~\Cref{sec:app-expt-nlp:mlm}, which is in turn initialized using the public RoBERTa checkpoint.
For the large size model, we lift weight from the state-of-the-art Pegasus model~\citep{zhang2019pegasus}, which is pretrained using an objective designed for summarization task.


To check if sparse attention causes significant degradation as compared to full attention, we further experiment on two shorter but popular datasets, where full attention can be used without significantly truncating the document. 
The statistics of these two datasets are in~\Cref{tab:sum_data}.
We see that our performance is competitive, which shows that sparse attention can achieve similar performance to a full attention models. 

\begin{table}
    \small
    \centering
    \begin{tabular}{@{}lrrrrrrrrrr@{}}
    \toprule
    & & \multicolumn{3}{c}{Instances} & & \multicolumn{2}{c}{Input Length}& & \multicolumn{2}{c}{Output  Length} \\
    \cmidrule{3-5} \cmidrule{7-8} \cmidrule{10-11}
    Dataset & & Training & Dev & Test & & Median & 90\%-ile & &  Median & 90\%-ile \\
    \midrule
       BBC XSum \citep{narayan2018don}  & &     204044 & 11332 & 11334 & &  359 & 920 & & 25 & 32 \\
       CNN/DailyMail  \citep{hermann2015teaching} & &  287113 & 13368 & 11490 & &   777 & 1439 & & 59 & 93 \\
    \bottomrule
    \end{tabular}
    \vspace{2mm}
    \caption{Shorter summarization dataset statistics.}
    \label{tab:sum_data}
\end{table}

\begin{table}
    \centering
    \small
    \begin{tabular}{@{}llrrrrrrrr@{}}
    \toprule
     \multicolumn{2}{l}{\multirow[b]{2}{*}{\hspace{-2mm}\normalsize{Model}}} & & 
     \multicolumn{3}{c}{BBC XSum} & & \multicolumn{3}{c}{CNN/DailyMail}\\
    \cmidrule{4-6} \cmidrule{8-10}
    & & & R-1 & R-2 & R-L & & R1 & R2 & R-L \\
    \midrule
    \multirow{9}{*}{\rotatebox[origin=c]{90}{Prior Art}}
    & Lead & & $16.30$ & $1.61$ & $11.95$ & & $39.60$ & $ 17.70$ & $ 36.20$ \\
    & PtGen~\citep{see2017get} & & $29.70$ & $ 9.21$ & $23.24$ & & $39.53$ & $17.28$ & $36.38$ \\
    & ConvS2S~\citep{gehring2017convolutional} & & $31.89$ & $11.54$ & $25.75$ & & $-$ & $-$ & $-$ \\
    & MMN~\citep{kim2018abstractive} & & $32.00$ & $12.10$ & $26.00$ & & $-$ & $-$ & $-$ \\
    & Bottom-Up~\citep{gehrmann2018bottom} & & $-$ & $-$ & $-$  & & $41.22$ & $18.68$ & $38.34$ \\
    & TransLM~\citep{khandelwal2019sample}  & & $-$ & $-$ & $-$  & & $39.65$ & $17.74$ & $ 36.85$ \\
    & UniLM~\citep{dong2019unified}  & & $-$ & $-$ & $-$  & & $43.47$ & $20.30$ & $40.63$ \\
    & Extr-Abst-BERT~\citep{liu2019text} & & 38.81  & 16.50  & 31.27 & &  42.13     & 19.60 & 39.18               \\
    & BART~\citep{lewis2019bart}  & & 45.14  & 22.27 & 37.25 & & 44.16 & 21.28 & 40.90 \\
    \midrule
    \multirow{4}{*}{\rotatebox[origin=c]{90}{Base}}
    & Transformer~\citep{vaswani2017attention}                & & 29.61 &  9.47 & 23.17 & & 34.89 & 13.13 & 32.12 \\
    & \; + RoBERTa~\citep{rothe2019leveraging} & & \underline{39.92} & \underline{17.33} & \underline{32.63} & & 39.44 & 18.69 & 36.80 \\
    & \; + Pegasus~\citep{zhang2019pegasus} & & 39.79 & 16.58 & 31.70 & & \underline{41.79} & \underline{18.81} & \underline{38.93} \\  
    & \bigb-RoBERTa & & 39.52 & 17.22 & 32.30 & & 39.25 & 18.46 & 36.61   \\
    \midrule
    \multirow{3}{*}{\rotatebox[origin=c]{90}{Large}}
    & Pegasus (Reported)~\citep{zhang2019pegasus} & & 47.60 & 24.83 & 39.64 & & 44.16 & 21.56 & 41.30 \\
    & Pegasus (Re-eval) & & \textbf{47.37} & \textbf{24.31} & \textbf{39.23} & & \textbf{44.15} & \textbf{21.56} & \textbf{41.05} \\
    & \bigb-Pegasus  & & 47.12 & 24.05 & 38.80 & & 43.84 & 21.11 & 40.74  \\
     \bottomrule
    \end{tabular}
    \vspace{2mm}
    \caption{Summarization ROUGE score for shorter documents.}
    \label{tab:app_sum_num}
\end{table}
